% SPDX-License-Identifier: Apache-2.0
\section{Exact-arithmetic specification and correctness}
\label{app:exact}

This is a first step towards formal verification: a precise specification
and a written mathematical proof of an exact-arithmetic model of
\code{walk()}. It is \textbf{not yet a machine-checked proof}, a proof about
the actual C++ floating-point execution, or a verification of the wrappers.
The correspondence below was checked against
\href{https://github.com/billlion/jacobsrays/blob/e1c4313550f86fd49547e7a9cb85f353d04f8a1d/include/jacobsrays/traverse.hpp}{\code{traverse.hpp} at revision \code{e1c4313}}.

A \emph{specification} states what a function must deliver, independently
of its implementation. A \emph{precondition} is an assumption on its inputs;
a \emph{postcondition} is a promised property of its output. A \emph{loop
invariant} is a statement that holds whenever execution reaches the top of
the loop. Proving an invariant requires establishing it initially and showing
that every iteration preserves it. Termination needs a separate argument.

\subsection{Inputs and the meaning of an output record}
\label{app:exact-contract}
Let $D\in\{2,3\}$. The preconditions are positive integers $n_a$, positive
real spacings $h_a$, and real origins $o_a$ and endpoints $p,e\in\mathbb R^D$,
where $0\leq a<D$. Define
\[
 r(t)=p+t(e-p),\qquad L=\|e-p\|_2,\qquad
 q_a=\frac{p_a-o_a}{h_a},\qquad v_a=\frac{e_a-p_a}{h_a}.
\]
The C++ core uses three coordinates; the 2D case embeds in it by adding a
stationary coordinate inside a single-cell slab.
Thus $q_a+t v_a$ is the ray coordinate in voxel units; $L$ is still a
physical length. The grid and its voxels are
\[
 \Omega=\prod_{a=0}^{D-1}[o_a,o_a+n_a h_a),\qquad
 V_c=\prod_{a=0}^{D-1}[o_a+c_a h_a,o_a+(c_a+1)h_a),
 \quad 0\leq c_a<n_a.
\]
These half-open boxes partition $\Omega$. In particular, a ray lying along
an internal face belongs to the higher-index cell in that coordinate.
The lower outer face is included and the upper outer face is excluded.

An output record $(c,a,b,\ell)$ denotes the portion with $a<t<b$ in voxel
$V_c$, with physical length $\ell=(b-a)L$. The implementation emits the
linear index $c_0+n_0(c_1+n_1c_2)$ in 3D, or $c_0+n_0c_1$ in 2D.
\textbf{The record need not own either endpoint as a point of its half-open
voxel.} For example, at $x=k$ a ray travelling towards decreasing $x$
enters cell $k-1$ immediately afterwards, although the point $x=k$ belongs
to cell $k$. Isolated endpoints have zero length.

The required postconditions are:
\begin{enumerate}
 \item If $L=0$, the output is empty. Otherwise, emit exactly one record for
 each voxel whose ray intersection has positive length, and no others.
 \item Each record has $0\leq a<b\leq1$, a valid cell index, and
 $r(t)\in V_c$ for every $a<t<b$. Its interval is the full positive-length
 intersection with that voxel, up to endpoints.
 \item Records occur in increasing parameter order. Consecutive intervals
 meet at an endpoint; their interiors are disjoint. Together they cover
 the part of the ray in $\Omega$, apart from finitely many endpoints.
 \item Each length is exact. In particular, writing $|S|$ for the
 one-dimensional measure of a parameter set,
 \[
   \ell_c=L\,|\{t\in[0,1]:r(t)\in V_c\}|,\qquad
   \sum_c\ell_c=L\,|\{t\in[0,1]:r(t)\in\Omega\}|.
 \]
 \item The algorithm terminates, with at most
 $1+\sum_a(n_a-1)$ emitted records when the output is nonempty.
\end{enumerate}
\subsection{The exact algorithm corresponding to the header}
\label{app:exact-algorithm}
Replace floating-point arithmetic by exact real arithmetic and fixed-width
indices by mathematical integers. Interpret \code{fma(first,v,q)} as
$q+\mathrm{first}\,v$, and the nested \code{hypot} calls as $L$.
The inactive-axis infinity is an ordered sentinel, used only in comparisons.
These replacements define the model; their correctness for machine
arithmetic is \emph{not} assumed to have been proved.
The callback is modelled as appending a record to a list. Any later proof
using an actual callback must specify its termination and effects on inputs.

Return immediately if $L=0$. For an axis with $v_a=0$, return empty if
$q_a\notin[0,n_a)$; otherwise it imposes no parameter restriction. For each
moving axis, form the closed slab interval with endpoints
$-q_a/v_a$ and $(n_a-q_a)/v_a$. Intersect these intervals with $[0,1]$,
calling the resulting lower and upper bounds $\alpha$ and $\beta$.
Return empty if $\alpha\geq\beta$.

Set $t=\alpha$, $x_a=q_a+\alpha v_a$, and initialise
\[
 c_a=\begin{cases}
       \lfloor x_a\rfloor,&v_a\geq0,\\
       \lceil x_a\rceil-1,&v_a<0,
     \end{cases}
 \qquad \sigma_a=\operatorname{sgn}(v_a).
\]
The next crossing parameter is
\[
 \tau_a=\begin{cases}
       (c_a+1-q_a)/v_a,&v_a>0,\\
       (c_a-q_a)/v_a,&v_a<0,\\
       +\infty,&v_a=0.
     \end{cases}
\]
While $t<\beta$, set $s=\min(\beta,\tau_0,\ldots,\tau_{D-1})$,
emit $(c,t,s,(s-t)L)$, and set $t=s$. If $s=\beta$, finish.
Otherwise, for \emph{every} axis with $\tau_a=s$, replace $c_a$ by
$c_a+\sigma_a$ and recompute $\tau_a$ using the formula above.

In the exact model, the header's clamps do not change the initial cell,
\code{stop > t} always holds in an entered loop iteration, and its early
out-of-grid return cannot occur before the final endpoint. The proof below
establishes these facts; they are not extra preconditions.

\subsection{Correctness proof}
\label{app:exact-proof}
\textbf{Theorem.} Under the preconditions above, the exact algorithm
satisfies all five postconditions.

\textbf{1. Clipping.}
If $v_a=0$, membership in that coordinate slab is constant and is exactly
the test $0\leq q_a<n_a$. If $v_a\ne0$, the inverse image of the closed
slab $[0,n_a]$ is the interval computed above. Replacing the actual half-open
slab by the closed slab changes membership only at its two endpoint
parameters. Hence, after the stationary-axis tests, the parameter set inside
the grid differs from the intersection of these closed intervals with
$[0,1]$ by at most finitely many points. It has positive length exactly when
$\alpha<\beta$. Every parameter strictly between these bounds lies in
$\Omega$. The empty returns, including point tangencies, are therefore
correct. The separately handled $L=0$ case has zero physical length.

\textbf{2. Initialisation.}
Assume now that $L>0$ and $\alpha<\beta$.
For $v_a>0$, clipping gives $0\leq x_a<n_a$: equality at the upper end
would leave no interior interval. Thus $\lfloor x_a\rfloor$ is a valid
cell, occupied immediately after $\alpha$. For $v_a<0$, similarly
$0<x_a\leq n_a$, and $\lceil x_a\rceil-1$ is the valid cell occupied
immediately after $\alpha$. For $v_a=0$, the stationary-axis test gives
$0\leq x_a<n_a$, and $\lfloor x_a\rfloor$ gives the prescribed ownership.
Both initial clamps in the header are consequently inactive.
For a moving axis, $\tau_a$ is strictly larger than $\alpha$ and is its
first cell-face crossing after $\alpha$.

\Needspace{12\baselineskip}
\textbf{3. Invariant.}
At the top of every entered iteration:
\begin{enumerate}[label=(\roman*)]
 \item $\alpha\leq t<\beta$, and the records already emitted correctly
 cover the parameter interval from $\alpha$ to $t$, up to endpoints.
 \item $c$ is a valid cell occupied immediately after $t$.
 \item Each moving axis has $\tau_a>t$, its first forward-in-parameter
 face crossing of this cell; stationary axes have $\tau_a=+\infty$.
\end{enumerate}
Here ``immediately after'' means that some $\varepsilon$ with
$0<\varepsilon\leq\beta-t$ exists such that $r(u)\in V_c$ for every
$u\in(t,t+\varepsilon)$. Initialisation proves the invariant before
the first iteration, with an empty output list.

\textbf{4. Preservation and the emitted interval.}
By (iii), $s>t$. No coordinate crosses a cell face for $t<u<s$, so the
entire open interval lies in $V_c$. Its physical length is exactly
$(s-t)L$. Appending it preserves ordered coverage and creates no overlap
of interiors. This is the maximal remaining interval in that cell:
at $s$ either the clipped ray ends or at least one coordinate leaves the cell.

If $s<\beta$, at least one moving axis has $\tau_a=s$.
Each such crossing is of an \emph{internal} plane: an outward crossing of
the outer box would end the clipped interval at or before $s$, contradicting
$s<\beta$. Updating every tied axis therefore gives valid indices and the
cell occupied immediately after $s$. A tied axis's next crossing increases
by $1/|v_a|>0$; untied axes already have $\tau_a>s$.
Thus (i)--(iii) hold for the next iteration. In particular, the header's
\code{next[a] <= stop} selects exactly the tied axes, and the early
out-of-grid return is unreachable here. Updating tied axes one at a time
within the update loop is harmless: no record is emitted between updates.

\textbf{5. Termination.}
The nonnegative integer
\[
 R(c)=\sum_{a:v_a>0}(n_a-1-c_a)+\sum_{a:v_a<0}c_a
\]
counts the internal planes still ahead, whether or not the finite ray will
reach them. Every nonfinal iteration decreases $R$ by the number of tied
axes, hence by at least one. This is a \emph{loop variant}: a quantity in
the nonnegative integers that strictly decreases on each continuation.
There cannot be infinitely many continuations. Since initially
$R\leq\sum_a(n_a-1)$, there are at most that many nonfinal iterations
and one final emission. The last emission ends at $\beta$.

\textbf{6. Completeness and uniqueness.}
The invariant and termination give contiguous coverage from $\alpha$ to
$\beta$, up to endpoints. Each moving cell coordinate only changes in its
fixed direction, and some coordinate changes between consecutive records;
therefore no cell can be visited twice. Every positive-length voxel
intersection contains an open parameter interval and so must be represented
in this coverage. Consequently the emitted records are exactly the required
intersections. Summing their lengths gives $L(\beta-\alpha)$, which equals
the length in the grid by the clipping argument. Valid cell coordinates
also imply a linear index between zero and $\prod_a n_a-1$ in mathematical
integer arithmetic. This proves all the postconditions.\hfill$\square$

\Needspace{18\baselineskip}
\subsection{A small example with a simultaneous crossing}
Take a $3\times2$ unit grid at the origin and
$p=(-1,-\tfrac12)$, $e=(4,2)$. Then
\[
 r(t)=(-1+5t,-\tfrac12+\tfrac52 t),\quad
 \alpha=\tfrac15,\quad\beta=\tfrac45,\quad L=\tfrac52\sqrt5.
\]
The exact output is
\begin{center}
\begin{tabular}{ccccc}
\toprule
Cell $c$ & Linear index & Entry $a$ & Exit $b$ & Length\\
\midrule
$(0,0)$ & $0$ & $1/5$ & $2/5$ & $\sqrt5/2$\\
$(1,0)$ & $1$ & $2/5$ & $3/5$ & $\sqrt5/2$\\
$(2,1)$ & $5$ & $3/5$ & $4/5$ & $\sqrt5/2$\\
\bottomrule
\end{tabular}
\end{center}
At $t=3/5$ the ray crosses $x=2$ and $y=1$ simultaneously, so both cell
coordinates advance. Cells $(2,0)$ and $(1,1)$ receive no positive-length
interval. The initial variant is $R=2+1=3$; the first crossing reduces it
to $2$ and the simultaneous crossing reduces it to $0$. The next record
ends at the grid boundary. Reversing the ray reverses the cell list and
maps every interval $(a,b)$ to $(1-b,1-a)$, preserving its length.

\subsection{What remains to be formalised}
The first Lean milestone would be the one-coordinate initialisation and
next-crossing lemmas, followed by the invariant and the decreasing integer
$R$. These statements involve exact arithmetic and order, not a numerical
error analysis. For rational input data, all clipping and crossing parameters
are rational; the generally irrational common factor $L$ can be handled
separately. A Lean model must still be linked to the stated algorithm, and
its proof must contain no unproved placeholders before being called checked.

The current proof does not cover rounding, near ties, overflow/underflow,
the C++ validators or casts, library-function accuracy, callback behaviour,
parallel accumulation, or Python/MATLAB memory handling. The output
cardinality bound and index bound above are over mathematical integers;
fixed-width implementations require further proofs. Comparing a reference
model with C++ would add useful tests, but would not establish their
equivalence for all inputs. That equivalence and a floating-point contract
are later milestones, separate from the exact theorem proved here.
